\section{Descriptive Statistics}

The descriptive section is organized around the empirical facts that motivate the regression analysis. We first document where workers are located in the firm-size distribution and how wages per hour vary across that distribution. We then
place the raw firm-size wage gap next to other remuneration gaps, describe worker composition by firm size, and examine whether the raw firm-size pattern appears within formality status, sex, education groups, and economic sectors.

Table~\ref{tab:descriptive-core-size-2008-2025} establishes the two core raw facts behind the paper. First, mean real hourly income rises sharply with firm size. In 2008, solo workers earned 5,388 COL 2025 pesos per hour on average, compared with 12,943 pesos among workers in firms with 101 or more employees, a ratio of 2.4 to 1. By 2025, the corresponding means were 6,534 and 16,859 pesos, a ratio of 2.6 to 1. Appendix Tables~\ref{tab:descriptive-core-size-2008-2025-101-reference} and \ref{tab:descriptive-core-size-2008-2025-11-19-reference} report the same means using alternative reference categories. Relative to 101+ workers, solo workers earned 42 percent of the large-firm mean in 2008 and 39 percent in 2025. Relative to workers in firms with 11--19 workers, 101+ workers earned 1.49 times as much in 2008 and 1.45 times as much in 2025. The all-worker mean increased from 7,752 to 10,295 pesos, equivalent to 1.7 percent annualized real growth between 2008 and 2025. Average weekly hours in the analytic sample declined from 47.8 to 43.8 over the same period, so real hourly wages grew faster than implied labor income per worker, which increased by 1.5 percent per year. Second, employment remains concentrated at the bottom of the firm-size distribution. In 2025, solo workers accounted for 35 percent of employment, and workers in firms with up to ten workers accounted for 61 percent. The middle of the distribution was thin: firms with 11--100 workers jointly accounted for 14 percent of employment. The 101+ category grew from 19 percent of employment in 2008 to 25 percent in 2025, but this increase did not eliminate the large employment mass in solo work and micro-firms.

\input{Paper/sections/descriptive_core_size_2008_2025_table}

This shift toward the largest category is relevant beyond the wage-premium
evidence. Recent growth theory emphasizes that the right tail of the firm-size
distribution tends to thicken with economic development, and that large firms
can matter for aggregate growth because they are central to idea search,
diffusion, and productivity dynamics \citep{Chen2026}. In this light,
Colombia's increase in the employment share of 101+ firms is consistent with
some movement toward a more developed firm-size structure. At the same time, the
continued concentration of employment in solo work and micro-firms shows that
this transformation remains incomplete.

Table~\ref{tab:descriptive-wage-gaps-2025} puts the raw firm-size gap in
context by comparing it with remuneration gaps across sex, formality status,
education, and economic activities in 2025. The wage gap by sex is small in
these weighted means: men earn 0.6 percent below the overall mean and women
0.8 percent above it. Other gaps are much larger. Formal workers earn 48.0
percent above the overall mean, while informal workers earn 39.1 percent below
it; this contrast is close in magnitude to the 2.6-to-1 ratio between workers
in 101+ firms and solo workers reported in
Table~\ref{tab:descriptive-core-size-2008-2025}. Education and sector
differences are larger still. Workers with higher education earn 71.6 percent
above the overall mean, while workers with preschool or less earn 57.9 percent
below it. Across activities, agriculture is 41.4 percent below the mean,
whereas information and communication is 112.8 percent above it. The raw
firm-size gap is therefore far larger than the wage gap by sex, comparable to
the formality gap, and smaller than the largest education and sector gaps.

\begin{table}[!htbp]
  \centering
  \caption{Mean real hourly labor income gaps by worker group, 2025}
  \label{tab:descriptive-wage-gaps-2025}
  \scriptsize
  \begin{tabular}{lp{0.34\textwidth}rrrr}
    \toprule
    Dimension & Group & Employment share & Mean income & Gap & Ratio \\
    \midrule
    Sex & Men & 58.8\% & 10,236 & -0.6\% & 1.00 \\
        & Women & 41.2\% & 10,379 & 0.8\% & 1.01 \\
    \midrule
    Formality & Formal & 44.9\% & 15,241 & 48.0\% & 2.43 \\
              & Informal & 55.1\% & 6,271 & -39.1\% & 1.00 \\
    \midrule
    Education & Preschool or less & 2.9\% & 4,333 & -57.9\% & 1.00 \\
             & Primary & 18.1\% & 5,693 & -44.7\% & 1.31 \\
             & Secondary & 45.5\% & 7,085 & -31.2\% & 1.63 \\
             & Higher education & 33.5\% & 17,667 & 71.6\% & 4.08 \\
    \midrule
    Activity & A - Agriculture & 13.8\% & 6,036 & -41.4\% & 1.00 \\
           & T - Household employers & 3.1\% & 6,454 & -37.3\% & 1.07 \\
           & I - Accommodation and food & 7.6\% & 6,930 & -32.7\% & 1.15 \\
           & H - Transport & 7.8\% & 7,775 & -24.5\% & 1.29 \\
           & F - Construction & 6.9\% & 8,406 & -18.3\% & 1.39 \\
           & G - Commerce & 17.0\% & 8,535 & -17.1\% & 1.41 \\
           & R/S - Arts and other services & 5.4\% & 8,765 & -14.9\% & 1.45 \\
           & C - Manufacturing & 11.0\% & 9,591 & -6.8\% & 1.59 \\
           & D/E - Utilities and water & 1.3\% & 11,295 & 9.7\% & 1.87 \\
           & B - Mining & 1.3\% & 11,441 & 11.1\% & 1.90 \\
           & L/M/N - Business services & 9.0\% & 13,459 & 30.7\% & 2.23 \\
           & Q - Health & 4.4\% & 14,859 & 44.3\% & 2.46 \\
           & P - Education & 4.2\% & 19,980 & 94.1\% & 3.31 \\
           & K - Finance and insurance & 1.9\% & 21,389 & 107.8\% & 3.54 \\
           & O - Public administration & 3.6\% & 21,508 & 108.9\% & 3.56 \\
           & J - Information and communication & 1.7\% & 21,903 & 112.8\% & 3.63 \\
    \bottomrule
  \end{tabular}
  \vspace{0.3em}
  \begin{minipage}{0.95\textwidth}
  \footnotesize
  Notes: Statistics use 2025 observations and GEIH expansion weights. Employment shares are computed within each dimension. Mean income is real hourly labor income in constant 2025 pesos. The gap is the percent difference in each group's weighted mean relative to the overall 2025 weighted mean. The ratio divides each group's mean by the lowest mean within the same dimension. Formality rows exclude pensioned occupied workers. Sector rows omit extraterritorial organizations.
  \end{minipage}
\end{table}


Table~\ref{tab:descriptive-worker-characteristics-2025} shows why the raw firm-size wage premium cannot be interpreted as a pure size effect. Larger firms employ a very different workforce. In 2025, only 8.3 percent of solo workers were formal, compared with 46.6 percent in firms with 6--10 workers and 98.0 percent in firms with 101 or more workers. Higher education also rises strongly with firm size, from 16.7 percent among solo workers to 62.9 percent in the 101+ category. The share of women varies less systematically across firm-size groups, while solo workers are older on average. These differences motivate the comparisons that follow, which ask whether the firm-size premium is visible within sex, formality, education, and sector groups.

\begin{table}[htbp]
  \centering
  \caption{Worker characteristics by firm size, 2025}
  \label{tab:descriptive-worker-characteristics-2025}
  \small
  \begin{tabular}{lrrrr}
    \toprule
    Firm size & Women & Formal & Higher education & Mean age \\
    \midrule
    Solo & 41.3\% & 8.3\% & 16.7\% & 44.5 \\
    2-3 & 34.6\% & 13.0\% & 17.0\% & 41.3 \\
    4-5 & 36.5\% & 24.2\% & 20.2\% & 37.6 \\
    6-10 & 40.5\% & 46.6\% & 29.9\% & 37.0 \\
    11-19 & 40.4\% & 64.4\% & 38.8\% & 37.2 \\
    20-30 & 42.7\% & 81.2\% & 45.4\% & 37.0 \\
    31-50 & 41.2\% & 90.4\% & 51.0\% & 37.5 \\
    51-100 & 43.7\% & 94.3\% & 53.4\% & 37.5 \\
    101+ & 45.8\% & 98.0\% & 62.9\% & 38.0 \\
    \bottomrule
  \end{tabular}
  \vspace{0.3em}
  \begin{minipage}{0.95\textwidth}
  \footnotesize
  Notes: Statistics use only 2025 observations and are weighted by GEIH expansion weights. Higher education corresponds to workers classified as superior or university educated.
  \end{minipage}
\end{table}


Figure~\ref{fig:income-formality-comparison} looks at the firm-size wage premium by formality status. The bubble size represents the percentage of formal or informal workers in each firm-size category in the corresponding year. Formal workers earn more than informal workers and are much more concentrated in large firms. The wage pattern by firm size is still visible within informal workers: among them, mean hourly income rises from 5,677 pesos for solo workers to 12,813 pesos in firms with 101 or more workers in 2025. Among formal workers, the raw profile is closer to a U-shape: formal solo workers earned 15,990 pesos per hour in 2025, formal workers in firms with 4--5 workers earned 12,036 pesos, and formal workers in firms with 101 or more workers earned 16,944 pesos. Table~\ref{tab:formal-worker-profile-2025} shows that formal solo work is a selected but heterogeneous category: 82.1 percent are own-account workers and 50.9 percent have higher education, with a sizable concentration in professional and administrative services, but 16.2 percent are domestic workers and 15.6 percent are in transport. Thus, the raw U-shape should be interpreted as evidence of worker composition within formal solo work, rather than as evidence that all formal solo workers are high-earning professionals.

\begin{figure}[!htbp]
  \centering
  \includegraphics[width=0.95\textwidth]{fig69_\figurelang.png}
  \caption{Mean real hourly income by firm size and formality status, 2008 and 2025}
  \label{fig:income-formality-comparison}
\end{figure}

\begin{table}[htbp]
  \centering
  \caption{Profile of formal workers by firm-size group, 2025}
  \label{tab:formal-worker-profile-2025}
  \small
  \begin{tabular}{lrrr}
    \toprule
    Characteristic & Formal solo & Formal 2--100 & Formal 101+ \\
    \midrule
    Share of formal workers & 6.4\% & 38.7\% & 54.9\% \\
    Mean real hourly income & 15,990 & 12,696 & 16,944 \\
    Mean age & 44.9 & 38.3 & 37.9 \\
    Women & 43.3\% & 41.9\% & 45.7\% \\
    Higher education & 50.9\% & 48.2\% & 63.4\% \\
    Own-account workers & 82.1\% & 7.7\% & 9.6\% \\
    Domestic workers & 16.2\% & 0.9\% & 0.0\% \\
    Professional and administrative services & 22.8\% & 14.6\% & 9.4\% \\
    Transport and storage & 15.6\% & 5.5\% & 5.6\% \\
    Commerce and repair & 14.7\% & 21.1\% & 10.1\% \\
    Health and social assistance & 5.2\% & 6.5\% & 10.6\% \\
    Households as employers & 16.2\% & 0.9\% & 0.0\% \\
    \bottomrule
  \end{tabular}
  \vspace{0.3em}
  \begin{minipage}{0.95\textwidth}
  \footnotesize
  Notes: Statistics use 2025 formal workers and GEIH expansion weights. The first row reports each column's share of all formal workers in the table. The 2--100 column pools all formal workers in firm-size categories from 2--3 through 51--100 workers. Higher education corresponds to workers classified as superior or university educated. Sector rows report the weighted share of each group located in the listed economic activity.
  \end{minipage}
\end{table}


Figure~\ref{fig:minwage-distribution-formality-size} shifts from mean wages to the distribution of hourly wages relative to the statutory hourly minimum in 2025. The vertical line marks one minimum wage, and the labels on the right report each firm-size category's employment share within the corresponding formality group. Among formal workers, the categories from 4--5 to 51--100 workers are visibly concentrated around one minimum wage. Appendix Table~\ref{tab:minimum-wage-bunching-2025} quantifies this pattern: between 36.9 and 42.7 percent of formal workers in these categories earn between 0.9 and 1.1 hourly minimum wages. Formal solo workers and formal workers in 101+ firms have wider upper tails, with 75th percentiles of 2.37 and 2.59 hourly minimum wages, respectively. Among informal workers, the distribution is shifted downward: 73.5 percent of solo workers and 65.9 percent of workers in firms with 2--3 workers earn less than 0.9 hourly minimum wages. This evidence helps interpret the flatter formal-sector profile: the minimum wage appears as a visible compression benchmark in formal wage setting, while informal employment contains much more mass below that benchmark.

\begin{figure}[!htbp]
  \centering
  \includegraphics[width=\textwidth]{fig77_\figurelang.png}
  \caption{Distribution of hourly wages relative to the statutory hourly minimum by firm size and formality status, 2025}
  \label{fig:minwage-distribution-formality-size}
\end{figure}

Figure~\ref{fig:income-sex-comparison} shows that the raw firm-size pattern is
similar for men and women. In 2025, among men, solo workers earned 6,816 pesos
per hour, compared with 16,929 pesos in firms with 101 or more workers; among
women, the corresponding means were 6,133 and 16,777 pesos. The large-firm
advantage is therefore very similar for both sexes. The figure also shows that wage
differences between men and women within a given firm-size category are small
relative to the difference between solo work and large-firm employment.

\begin{figure}[!htbp]
  \centering
  \includegraphics[width=0.95\textwidth]{fig70_\figurelang.png}
  \caption{Mean real hourly income by firm size and sex}
  \label{fig:income-sex-comparison}
\end{figure}

Figure~\ref{fig:income-education-comparison} shows that the raw firm-size
premium is visible within every education group. Among workers with preschool
or less, mean hourly income in 101+ firms is 2.1 times solo-worker income. The
corresponding 101+-to-solo ratio is 1.8 among workers with primary education,
1.5 among workers with secondary education, and 1.7 among workers with higher
education. Thus, the large-firm advantage is not only a comparison between more
and less educated workers; it is also visible within each education level.

\begin{figure}[!htbp]
  \centering
  \includegraphics[width=\textwidth]{fig72_\figurelang.png}
  \caption{Mean real hourly income by firm size and education group}
  \label{fig:income-education-comparison}
\end{figure}

Figure~\ref{fig:income-sector-comparison} shows that the firm-size wage premium
is widespread across economic activities. Mean hourly income in firms with 101 or more workers exceeded
solo-worker income in 14 of 16 sectors in 2008 and in all 16 sectors in 2025.
The two 2008 exceptions were finance and insurance and business services, where
solo workers earned slightly more than workers in the largest firms; by 2025,
the large-firm category also paid more in both sectors. The size premium is
especially pronounced in utilities and water, households as employers, mining,
manufacturing, accommodation and food, and transport, while it is much flatter
in information and communication, health, finance, and business services. These
sector panels therefore reinforce the aggregate pattern while also showing that
the size premium varies substantially across activities.

\begin{figure}[p]
  \centering
  \includegraphics[width=\textwidth]{fig71_\figurelang.png}
  \caption{Mean real hourly income by firm size and sector, 2008 and 2025}
  \label{fig:income-sector-comparison}
\end{figure}

Taken together, the descriptive evidence motivates the regression analysis
below. The central pattern is not only that mean hourly income is higher in
larger firm-size categories. The tables and figures also provide suggestive
evidence that a firm-size wage premium appears within each sex, within each
education group, in most economic activities, and even within both formal and
informal employment. At the same time, the descriptive results show strong
sorting by formality, education, sector, and position in the firm-size
distribution. The econometric exercise therefore asks how much of the raw
firm-size premium survives after controlling for observable worker
characteristics, formality, and sector-year conditions.
